\documentclass[11pt]{article}
\usepackage[margin=26mm]{geometry}
\usepackage{amsmath,amssymb,amsthm,booktabs,longtable,array,hyperref}
\hypersetup{colorlinks=true,linkcolor=blue,urlcolor=blue}
\newtheorem{theorem}{Theorem}
\newtheorem{lemma}{Lemma}
\newcommand{\CE}{C_{\mathrm E}}
\newcommand{\pos}[1]{(#1)_+}
\newcommand{\R}{\mathbb R}
\title{Three further reductions for the conditional LOCAL42 route}
\author{Research and independent audit in the user's fork}
\date{27 October 2026}
\begin{document}
\maketitle
\begin{abstract}
We give an elementary replacement for all 65 zero-loss support calls, an
exact three-candidate formula for the LOW curvature supremum, and a tighter
resource argument that removes two selected LOW producers. These are
conditional components, not a proof of the sharp iid Berry--Esseen
conjecture. The full acquisition, equality-compatible attainment, contact,
local-comparison and finite-base obligations remain. No original certificate,
target or production checker is changed. A separate one-pass execution
prototype is discussed only at its tested protocol scope.
\end{abstract}

\section{Scope and preserved premises}
The user's fork is \path{kolodziejek/berry-esseen-certificates}.
The immutable reference and this continuation's starting commit are,
respectively:
\begin{center}\small\ttfamily
6346fc1b4432bc5982822aa696eb566ed9e1e1ee\\
42433b448d9389359e642a3384a949ffdd87b903
\end{center}
Exact source paths, hashes and literal line locations are in the companion
notes:
\begin{quote}\ttfamily
elementary-support.md\\
local-kernel.md\\
tighter-resources.md
\end{quote}
Their small checkers read the pinned source
archive without executing its native arithmetic kernels.

The earlier reduction retains a genuine sharp base through 41. At a least
index $n\ge42$ with $C_n\ge\CE$, both required predecessors are sharp,
the global predecessor decrement is zero, and support-free smoothing gives
\[
 n(C_n-\CE)<\frac{3003}{10250}<0.3,\qquad
 C_n<\frac{73}{175}<0.4226.
\]
Both deleted-score caps are below $0.4253$. Those scalar conclusions alone
do not establish LOCAL42. In particular they do not imply zero conditional
noise, compact support, attainment, or a numerical certificate.

The support result below applies to an attained positive closed maximizing
pair $(P,H)$ with standardized iid marginal, third absolute moment
$\rho=E|X|^3$, and signed moment $\beta=E[X|X|]$. All moment and same-law
contact premises are retained. The resource result is applied only
\emph{after} the corresponding confinement and acquisition. Its extension
from the literal finite $42$--$48$ source scope requires the separate
uniform42 bridge; the floor $42$ alone does not prove that extension.

\section{Elementary support without a positivity mesh}
Throughout this section, the actual integer satisfies $n\ge42$ and the
attained score satisfies $C_n\ge\CE$.
Set $v=n-1$, $s=\sqrt{n/v}$ and $F_v(x)=\Phi(x/\sqrt v)$. The retained
contact inequality on the complete closed support is
\begin{equation}\label{contact}
 0\le H_n(H,t)+C_n\{-|t|^3+\tfrac32\rho t^2+3\beta t+\rho e_n\},
\end{equation}
where
\[
 e_n=n\bigl(\sqrt{n/(n-1)}-1\bigr)-\tfrac12<\frac{0.384}{n},
\]
and the exact Gaussian expression is
\[
 \frac{H_n(H,t)}{n^{3/2}}
 =F_v(H-t)-F_{v+1}(H)+tF_{v+1}'(H)
       -\frac{t^2-1}{2}F_{v+1}''(H).
\]
The same-owner predecessor score is at most the global predecessor score;
at zero global decrement the discarded signed-loss term is nonpositive.

\subsection{Fourth-order expansion with uniform remainder}
Using $\partial_vF_v=F_v''/2$, expand $F_v(H-t)$ spatially through order
four, $F_{v+1}$ through second order in variance, and its first and second
spatial derivatives through first order in variance. The order-zero, linear
and quadratic spatial terms cancel. The remaining cubic and quartic
coefficients are $(-t^3+3t)/6$ and $(t^4-6t^2+3)/24$. Thus, for $t=\pm u$,
\begin{equation}\label{expansion}
 \frac{H_n(H,\pm u)}{s^3}
  =\pm A(1-z^2)\phi(z)+B\phi'''(z)+R,\qquad z=H/\sqrt v,
\end{equation}
where
\[
 A=\frac{u^3-3u}{6},\qquad B=\frac{u^4-6u^2+3}{24\sqrt v}.
\]
For $u\ge5/2$, both are positive. The four Taylor remainders, after division
by $s^3$, have bounds
\[
 \frac{u^5\|\phi^{(4)}\|_\infty}{120v},\quad
 \frac{\|\phi^{(5)}\|_\infty}{48v^{3/2}},\quad
 \frac{u\|\phi^{(4)}\|_\infty}{8v},\quad
 \frac{(u^2-1)\|\phi^{(5)}\|_\infty}{16v^{3/2}}.
\]
Gaussian Fourier inversion gives
$\|\phi^{(4)}\|_\infty\le3\phi(0)<6/5$ and
$\|\phi^{(5)}\|_\infty\le8/\pi<51/20$. Therefore
\begin{equation}\label{error}
 |R|\le E_v(u):=\frac{u^5/100+3u/20}{v}
       +\frac{51(3u^2-2)}{960v^{3/2}}.
\end{equation}
Every bound is uniform in the raw threshold $H$. No seventh derivative
of $F_v$, sixth derivative of $\phi$, or signed fifth-order tangent is used.
The rational lower bound $\pi>281476/89625>160/51$ suffices for these
constants and $\phi(0)<2/5$.

\begin{lemma}[Two all-threshold envelopes]\label{envelopes}
For $u\ge5/2$,
\[
 H_n(H,\pm u)\le s^3\{\phi(0)(A+2B)+E_v(u)\}.
\]
If also $B\le A/4$, then
\[
 H_n(H,\pm u)\le
 s^3\left\{\phi(0)\left(A+\frac{9B^2}{4A}\right)+E_v(u)\right\}.
\]
Both assertions hold for every real $H$ and both signs separately.
\end{lemma}
\begin{proof}
The first follows from
$\|(1-z^2)\phi(z)\|_\infty=\phi(0)$ and
$\|\phi'''\|_\infty<2\phi(0)$.
For the positive sign and $|z|\le1$, the retained terms in
\eqref{expansion} are at most
$\phi(0)(A-Az^2+3B|z|)$; completing the square proves the sharper bound.
For $|z|>1$ the first term is nonpositive and the second is at most
$2B\phi(0)\le A\phi(0)/2$.
For the negative sign, its first term is at most
$2e^{-3/2}A\phi(0)<A\phi(0)/2$, and the second is at most
$2B\phi(0)\le A\phi(0)/2$. Add the uniform remainder.
\end{proof}

\subsection{Rational boundary inequalities and full rays}
Use
\[
 S=81/80,\quad \ell=6,\quad c_0=819/2000<\CE,\quad
 H_4=u^4-6u^2+3,
\]
so $s\le S$, $\sqrt v>\ell$ and $e_n<8/875$. Define
\begin{align*}
 E(u)&=\frac{u^5}{4100}+\frac{3u}{820}
                   +\frac{17(3u^2-2)}{78720},\\
 Q_{r,b}(u)&=-u^3+\tfrac32ru^2+3bu+\tfrac8{875}r,\\
 P_{r,b}(u)&=\tfrac25S^3(A+H_4/72)+S^3E(u)+c_0Q_{r,b}(u),\\
 P^\sharp_{r,b}(u)&=\tfrac25S^3
       \left(A^2+\frac{H_4^2}{256\cdot41}\right)
       +S^3E(u)A+c_0Q_{r,b}(u)A.
\end{align*}
For acquired $\rho\le r$, $|\beta|\le b$, the contact coefficient is at
most $Q_{r,b}$. Check $Q_{r,b}<0$ before replacing its actual multiplier
$C_n$ by the lower constant $c_0$. Lemma~\ref{envelopes} then makes
$P<0$ sufficient at a boundary; $P^\sharp<0$ is sufficient when $B\le A/4$.

To extend a boundary to a whole ray, put
\[
 L_{r,b}(u)=\tfrac25(u^2/2+1/2+8/875)+c_0Q'_{r,b}(u).
\]
The retained direct derivative bound is
$|\partial_tH_n|\le(2/5)(t^2/2+1/2+e_n)$.
If $Q,Q',Q'',L,L',L''$ are negative at $U$, their required signs persist
for $u\ge U$, since $Q'''=-6$ and $L''=2/5-6c_0<0$.
The contact upper bound for each sign then decreases strictly along the
whole ray. This includes endpoints and atoms.

\paragraph{Initial case.}
Take $R=96/25$, $r=1841/1000$, $b=321/500$. The retained moment inequalities
give
\[
 \rho<
 \min\left\{\frac{.3031\cdot.646}{c_0-.3031},
              \frac{.3322\cdot.429}{c_0-.3322}\right\}
 =\frac{7361}{4000}<r.
\]
The sign-moment inequality gives $|\beta|<b$, using
\[
 (1+b^2)^2-r^2(1-b^2)=\frac{93288849}{50000000000}>0.
\]
At $u=R<6$,
$H_4=6uA-3(u^2-1)<36A$ and $B<H_4/144<A/4$. Exact arithmetic gives
\[
 -P^\sharp_{r,b}(R)=
 \frac{2214258283536204102723}{10250000000000000000000}>\frac15.
\]
All six ray signs are negative, proving support in $(-R,R)$.

\paragraph{Common LEFT bridge.}
Take $R=121/40$, $r=143/100$, $b=93/200$, outwardly enclosing the acquired
common moment maxima $71458741/50000000$ and $46459359/100000000$.
Then
\[
 -P_{r,b}(R)=\frac{2533499510174559}{214958080000000000}>\frac1{100}.
\]
All six ray signs are negative.

\paragraph{A continuous family for the other 63 calls.}
All remaining acquired parameter tuples satisfy
\[
 \frac52\le R\le\frac{111}{40},\qquad
 \rho\le r_R=\frac7{20}R+\frac13,\qquad
 |\beta|\le b_R=\frac38R-\frac35.
\]
Let $p(R)=-P_{r_R,b_R}(R)$ and $w=R-5/2$. The coefficients of
$p(R)=\sum_{j=0}^5a_jw^j$ are
\[
\begin{array}{c|r}
j&a_j\\ \hline
0&123579925563/1343488000000\\
1&-467155860219/1679360000000\\
2&705959271/13434880000\\
3&8703567/167936000\\
4&-7499223/839680000\\
5&-531441/2099200000
\end{array}
\]
With $h=11/40$, their signs imply throughout $0\le w\le h$
\[
 p'(R)\le a_1+2a_2h+3a_3h^2
       =-\frac{797749885413}{3358720000000}<0.
\]
Consequently
\[
 p(R)\ge p(111/40)
       =\frac{4403736280322409}{214958080000000000}>\frac1{50}.
\]
For rays, freeze each pair $(r_R,b_R)$. The common larger caps
$r_*=131/100$, $b_*=441/1000$ have at $u=5/2$ the six signs
\[
 -\frac{16991}{700000},\quad-\frac{3801}{500},\quad-\frac{1107}{100},
 \quad-\frac{11615533}{7000000},\quad-\frac{706633}{200000},
 \quad-\frac{2057}{1000}.
\]
They majorize the corresponding signs and derivative bounds for every fixed
pair. Combining the negative boundary with these complete rays proves all
63 calls. Continuing the affine caps along an infinite ray would be invalid.

\begin{theorem}[Conditional elementary support replacement]
For an actual integer $n\ge42$ and attained score $C_n\ge\CE$, under the
retained contact, zero-loss and ordered same-law acquisition
hypotheses, all 65 support conclusions used by the zero-loss LOCAL42 route
follow from the preceding three cases. The old radial Bernstein cells,
the newer 14-group/672-coefficient support certificate, the signed
fifth-order identity, its fourth-derivative lower tangent and its associated
domination gates are unnecessary for this alternative support component.
\end{theorem}
The input bindings and eighteen common ray signs are checked by the script
\begin{center}\ttfamily check\_elementary\_support.py\end{center}
These arithmetic facts do not prove
the acquired moment inequalities themselves.

\section{Exact curvature maximization and safe transport}
For real $b$, $1/2<c\le1$ and $[l,h]\subseteq[-1/2,1/2]$, define
\[
 f(u)=\frac{\pos{b+u-1}}{(c-u)^2}
       +\frac{\pos{b+u}}{[c+\min(u,c/3)]^2},\qquad u_*=c-2b.
\]
All denominators are strictly positive. In the kernels $c=1-d_L$ and
$b=p_+-\min(H/\ell)$, where the minimum uses all four signed endpoint
quotients. The variable $b$ is not a probability and cannot be freely
reflected.

\begin{theorem}[Exact three-candidate supremum]
The supremum of $f$ over $[l,h]$ is the maximum of $f(l)$, $f(h)$, and
the value $1/[4(c-b)]$ if
\[
 c>b,\qquad l\le u_*\le h,\qquad
 u_*\le c/3,\qquad u_*\le1-b.
\]
An ineligible third value is omitted. The formula includes singleton
intervals, coincident kinks and $c=1$.
\end{theorem}
\begin{proof}
Write $f=r+g$ for the two summands. On $b+u>1$,
\[
 r''=\frac{2[3(b-1)+u+2c]}{(c-u)^4}>0,
\]
since its bracket is $3(b+u-1)+2(c-u)>0$. Where $u<c/3$ and $b+u>0$,
\[
 g'=\frac{c-2b-u}{(c+u)^3},\qquad
 g''=\frac{2(3b+u-2c)}{(c+u)^4}.
\]
If also $r>0$, then $3b+u-2c>3-2u-2c\ge0$, so $f$ is convex there.
For $u>c/3$, $g$ is convex piecewise linear. At each of
$-b$, $1-b$ and $c/3$, the derivative has a nonnegative jump; simultaneous jumps
add. Thus none is a strict interior maximum.

Where $r=0$, either $g=0$, or below $c/3$ its derivative changes from
positive to negative only at $u_*=c-2b$. Its positive numerator condition
is $c>b$, and $r=0$ there is equivalent to $u_*\le1-b$.
Its value is $1/[4(c-b)]$. Above $c/3$, $g$ is nondecreasing.
These cases locate a global maximum among the stated candidates.
Continuity handles boundaries and collisions; constant segments contribute
a value already attained at an endpoint or a listed maximizer.
\end{proof}

For an old grid interval $[a,d]$ and every $u\in[a,d]$,
\[
 f(u)\le \frac{\pos{b+d-1}}{(c-d)^2}
       +\frac{\pos{b+d}}{[c+\min(a,c/3)]^2}.
\]
Therefore the exact maximum is at most the old 64-interval outward bound.
For fixed $b,u$, $f$ is nonincreasing in $c$. CENTRAL's first cached band
has the smallest $c$, so its cached envelope also dominates later exact
values; this requires the original band order and joint-gap relation.
LEFT retains its band-dependent cache.

\subsection{Why the final correction needs a zero clamp}
The curvature improvement can make the recomputed inverse-index correction
negative. One must not assume the old assertion $I>0$ remains true.
Within a fixed positive-noise cell, the actual forms are
\[
 B(\lambda)=z+\frac{q}{a\lambda+G^2}+s\lambda-t,\qquad
 I_{\rm new}(\lambda)=A_0+\frac{A_1}{\lambda},
 \quad q\ge0,\ a>0,\ A_1\ge0.
\]
Only the curvature contribution to $A_0$ decreases. Put
$I_+(\lambda)=\max(0,I_{\rm new}(\lambda))$. For $n\ge N$,
\[
 mB(\lambda)+\frac{I_{\rm new}(\lambda)}n
 \le mB(\lambda)+\frac{I_+(\lambda)}N.
\]
Both $B$ and the right-hand expression are convex for $\lambda>0$,
including a possible interior zero-clamp kink, since $m>0$.
Endpoint checks thus cover the closed interval. At an accepted old endpoint,
\[
 I_{\rm new}\le I_{\rm old},\ I_{\rm old}>0
 \quad\Longrightarrow\quad 0\le I_+\le I_{\rm old}.
\]
The original strict target $mB+I_{\rm old}/N<-1/1000$ is preserved.
The separate $B<0$ gate remains necessary to use the mass lower bound.

This is conditional transport from an accepted old cell, not a full replay
or a suppression of an old assertion. A new implementation replaces 64
candidate values by at most three per curvature cache miss. All marked-tail,
normalization, mass, joint-resource and other comparison work remains.

\section{Retaining the score-excess charge}
Let the actual integer satisfy $n\ge N=42$. Put $c=\CE$, with actual score $C\ge c$,
$\rho=E|X|^3\ge1$, and $\kappa=997/15000$. Retain the source phase resource
\begin{equation}\label{resource}
 n(C-c)\rho+\kappa nD+n^{3/2}L\le\mathcal B(H),\qquad L\ge0,
\end{equation}
with its original same-law geometry, support, signed phase, concentration,
moment and deleted-score hypotheses. In particular the conserved phase
curvature is not charged twice.

For each of two acquired moment corners $j$ and each phase sign, its
quadratic upper envelope is
\[
 A_j(C)+F_{j,\pm}(C)|H|-Q_j(C)|H|^2.
\]
Here $A_j,Q_j$ are affine in $C$, and each signed $F$ is a finite maximum
of affine functions plus a score-independent slope. On the endpoint-score
interval $[C_0,C_1]$, its chord $\widehat F$ majorizes $F$.
The supremum of these chord quadratics is
\[
 G(C)=\max_j\left\{A_j(C)+
 \frac{\max(0,\widehat F_{j,+}(C),\widehat F_{j,-}(C))^2}{4Q_j(C)}
 \right\}.
\]
Positive endpoint $Q_j$ gives positivity throughout. As a supremum of
affine functions of $C$ indexed by $j$, sign and $|H|\ge0$, $G$ is convex.

\begin{lemma}[Two-interval score charge]\label{charge}
Suppose $C_0<c<\chi<C_1$ and the actual score is at most $C_1$.
Then \eqref{resource} implies
\[
 \kappa nD\le
 B_*=\max\{G(C_0),G(\chi),G(C_1)-N(C_1-\chi)\}.
\]
\end{lemma}
\begin{proof}
For actual $C\in[c,\chi]$, discard its nonnegative excess charge and
bound $G(C)$ by the endpoints of the larger formal interval $[C_0,\chi]$.
For $C\in[\chi,C_1]$,
$n(C-c)\rho\ge N(C-\chi)$, using only $\rho\ge1$.
The convex function $G(C)-N(C-\chi)$ is bounded above by its endpoint
values. Take the maximum. This includes $C=c$ without assigning a
positive charge at equality.
\end{proof}

For LEFT, $C_0=.4095$, $C_1=.4226$ and $\chi=.40974$.
Exact arithmetic on all 45 rows gives $B_*=G(\chi)$.
For completeness, with
$s_k(x)=\sum_{i=0}^{k-1}(-1)^ix^{2i+1}/(2i+1)$, the alternating-series
bounds and Machin's tangent-addition identity give
\[
 16s_6(1/5)-4s_1(1/239)<\pi
 <16s_5(1/5)-4s_2(1/239).
\]
Together with $3.162277660<\sqrt{10}<3.162277661$ and
$c^2=(19+6\sqrt{10})/(72\pi)$ these prove $C_0<c<\chi$ using rationals.

\subsection{Keep one joint residual}
The old weighted bound is $nD\le Nd_{\rm old}$; define
\[
 d_{\rm new}=\min\{d_{\rm old},B_*/(\kappa N)\}.
\]
This is a bound on $nD$, not a bare bound on $D$ subsequently rescaled.
Retain the acquired geometry $\Theta,A,B$ and closed probability panel
$[l,u]$. Write $s_l=3.162277660$, $s_u=3.162277661$, and set
\begin{align*}
 v_*&=\max_{p\in[l,u]}p(1-p),\\
 \delta_*&=\max\{0,(4-s_u)/2-u,\ l-(4-s_l)/2\},\\
 b_*&=(6+2s_l)
   \frac{\sqrt{(1-\Theta)^3}_{\,\rm lower}}
        {\sqrt{v_*}_{\,\rm upper}}\,\delta_*^2,\\
 K&=\min\{3(s_l+1)B,\ 3(s_l+3)A\}>0,\qquad
 r=d_{\rm new}-b_* .
\end{align*}
The retained same-law inequality is
\[
 nD\ge Nb_*+6N(1-\Theta)(A+B)d_{\rm in}
                 +\frac{K(1-\Theta)}{v_*}\lambda.
\]
If $r<0$, the row is impossible. If $r>0$, put
\[
 \Lambda_{\rm new}=\frac{Nv_*r}{K(1-\Theta)},\qquad
 d_{{\rm new,in}}=\frac{r}{6(1-\Theta)(A+B)}.
\]
Then the \emph{single} residual gives
\[
 \frac{d_{\rm in}}{d_{{\rm new,in}}}
       +\frac{\lambda}{\Lambda_{\rm new}}\le1.
\]
It is not spent independently on noise and inward displacement. No tested
row has $r=0$; that case would require a separate zero-noise join.

CENTRAL exports upward-rounded $A,F$ and downward-rounded $Q$, and only an
outward raw score quantity. Its source rounding contract gives
$\mathrm{raw}_{\rm exact}\in[\mathrm{raw}_{\rm out}-2\cdot10^{-6},
\mathrm{raw}_{\rm out}]$. Accordingly its high score endpoint belongs to
\[
 [\,\min(.4247,\chi+(\mathrm{raw}_{\rm out}-2\cdot10^{-6})/42),
    \min(.4247,\chi+\mathrm{raw}_{\rm out}/42)\,].
\]
Enclose the chord weight over this interval; use upper $A,F$, lower $Q$,
and the lower high-score endpoint in the subtraction in
Lemma~\ref{charge}. This gives valid outward budgets for all 47 central rows.

\subsection{Exact removable sectors}
The following decimal ceilings are rounded upwards from exact fractions.
Rows are zero-based. The full 92-row scan establishes all newly crossed
thresholds among $.05,.08,.10,.12$.
\begin{center}
\begin{tabular}{@{}llp{69mm}@{}}
\toprule
Row & New conclusion & Consequence after acquisition\\
\midrule
LEFT 0 & $r<-.0003$ & Entire final row excluded; TINY comparison unnecessary.\\
LEFT 5 & $\lambda\le.049944478$ & MID domain above $.05$ empty.\\
LEFT 8 & $\lambda\le.077273055$ & Both signed LOW producers above $.08$ unnecessary.\\
LEFT 10 & $\lambda\le.094055578$ & Domain above $.10$ empty; two profile calls on $[.10,.12]$ unnecessary.\\
LEFT 13 & $\lambda\le.117287323$ & Domain above $.12$ empty; two profile calls on $[.12,.30]$ unnecessary.\\
CENTRAL 28 & $\lambda\le.099809552$ & Comparison sector above $.10$ empty; profile already inherited.\\
\bottomrule
\end{tabular}
\end{center}
Restriction of the old complete cover to these smaller same-law domains
preserves signs, nearest counts, closed phase seams and probability
endpoints. Each removed closed noise band has its lower endpoint strictly
above the new ceiling. The separate $\lambda=0$ Bernoulli theorem remains.
LEFT 6, 9 and 14 and CENTRAL 38 still exceed their relevant thresholds;
this method does not justify deleting them.

Thus a newly named upper-module contract could use 150 LOW producers
(78 CENTRAL, 72 LEFT), with 19 native obligations still retained. The
candidate task count becomes 169, before new joins. This is not a count
for the whole repository or a measured runtime saving. In particular,
row 0 is excluded only after its final resources have been acquired:
none of its acquisition or confinement ancestors is removed.

\section{One-pass protocol and compatibility boundary}
The old selected plan evaluates 152 LOW configurations in producers and
evaluates them again in binders. The separately named prototype
\texttt{onepass\_low.py} retains the original 152 configurations, source
pins, serialized rows, index hashes, coverage and nine binder controls per
family. It owns fresh evaluator outputs in memory, runs the original
non-replay binder validations, and compares read-back bytes and summaries
to those fresh outputs. It does not accept user-supplied rows or receipts
as a replacement for fresh evaluation.

Synthetic fixtures using the small original binder bodies produce identical
legacy rows, indices and summaries with 152 rather than 304 evaluator calls;
14 further rejection controls pass. The native source-loading and full
arithmetic path has not been executed, so this is protocol evidence, not
a production acceptance result.

The curvature formula and the tighter resources change numerical outputs
or the selected labels. They therefore do \emph{not} satisfy the prototype's
old output/index contract. Combining the three changes requires a new,
complete contract and independent review. Their potential savings must not
be multiplied into an unsupported overall speedup.

\section{Evidence and remaining obligations}
The component reports and \texttt{runs/reduction2-20261027/} separate
mathematical and execution status. The support checker supplies exact
three-case inequalities, all 65 source bindings and negative controls.
The curvature checker tests 3,648 calls on the actual manifest grids,
with 62,942 exact point diagnostics; these finite checks supplement the
continuous proof, not replace it. Its first workload count used defaults
and was corrected using actual producer settings, preserving the old run.
The resource checker reconstructs both old LEFT intercepts on all 45 rows,
then calculates 92 outward consequences with 12 rejection controls.

Independent reviews use separate worktrees and disclose construction roles.
Their reports contain independent support algebra, curvature adversarial
checks, alternative outward resource arithmetic, and the root's protocol
review. Review agreement is not an additional mathematical premise.

The genuine base provider, equality-compatible attainment, same-law
contact and moment analysis, strict confinement, Fourier bounds, remaining
local acceptances, zero-noise Bernoulli input, lower approximation and
complete composition contract remain obligations. No full LOCAL42,
whole-proof PASS, or proof of the conjecture is asserted here.
\end{document}
